\documentclass[11pt]{article}
\usepackage[margin=0.95in]{geometry}
\usepackage{amsmath,amssymb,amsthm}
\usepackage[T1]{fontenc}
\usepackage{lmodern}
\usepackage[hidelinks]{hyperref}
\usepackage{microtype}
\setlength{\parindent}{0pt}
\setlength{\parskip}{6pt}
\newtheorem{proposition}{Proposition}
\begin{document}
\begin{center}
{\Large A cubic obstruction to Conjecture 00000003315\par}
\vspace{8pt}
gaochengzhecpu \qquad October 4, 2026
\end{center}
\begin{abstract}
The toric ideal of the full lattice-point configuration of the triangle
$\operatorname{conv}\{(0,0),(2,1),(1,2)\}$ contains the nonzero cubic
$X_1X_2X_3-X_0^3$, but contains no nonzero binomial of degree at most two.
It therefore cannot be generated by quadratic circuit relations, as
asserted in Conjecture 00000003315. The accompanying Lean project proves
this for an actual polynomial algebra map, its kernel ideal, and every
putative set of quadratic binomial generators.
\end{abstract}

\section{The lattice polygon and its toric map}
Let
\[
 P=\operatorname{conv}\{(0,0),(2,1),(1,2)\}.
\]
Its complete set of lattice points, in the order used below, is
\[
 p_0=(1,1),\qquad p_1=(2,1),\qquad p_2=(1,2),\qquad p_3=(0,0).
\]
Indeed, a point of $P$ has the form $(x,y)=(2b+c,b+2c)$ with
$b,c\geq0$ and $b+c\leq1$. Thus
\[
 x,y\geq0,\qquad x+y\leq3,\qquad x\leq2y,\qquad y\leq2x.
\]
For integral $x,y$, either one coordinate is zero and both are zero,
or both are positive and $x+y\leq3$ gives exactly $(1,1),(2,1),(1,2)$.
Conversely the three vertices belong to $P$, and $(1,1)$ is their
barycenter. This proves the claimed complete lattice-point list.

Homogenize the points to the columns $a_i=(1,p_i)$. Over $\mathbb Q$,
consider the polynomial algebra homomorphism
\[
\begin{aligned}
 \varphi:\mathbb Q[X_0,X_1,X_2,X_3]&\longrightarrow\mathbb Q[t,x,y],\\
 (X_0,X_1,X_2,X_3)&\longmapsto(txy,\,tx^2y,\,txy^2,\,t).
\end{aligned}
\]
The toric ideal of this homogenized polytopal lattice configuration is
$I=\ker\varphi$. The homogeneous first coordinate is essential here:
in particular, $p_3=(0,0)$ maps to $t$, not to the constant $1$.

\begin{proposition}
$I$ is nonzero and has no generating set consisting of binomials of
degree at most two. In particular, it is not generated by quadratic
circuit relations.
\end{proposition}
\begin{proof}
The distinct monomials in
\[
 F=X_1X_2X_3-X_0^3
\]
have the same image $t^3x^3y^3$. Hence $0\ne F\in I$.

For $e=(e_0,e_1,e_2,e_3)\in\mathbb N^4$, the image of $X^e$ is the
monomial with exponent vector
\[
 A e=(e_0+e_1+e_2+e_3,\ e_0+2e_1+e_2,\ e_0+e_1+2e_2).
\]
Suppose $Ae=Af$ and $|e|,|f|\leq2$. Writing $d=e-f\in\mathbb Z^4$,
the three coordinate equations give
\[
 d=(-3r,r,r,r)\qquad\text{for some }r\in\mathbb Z.
\]
If $r>0$, then $f_0=e_0+3r\geq3$, contrary to $|f|\leq2$.
If $r<0$, then $e_0=f_0-3r\geq3$, contrary to $|e|\leq2$.
Thus $r=0$ and $e=f$.

Now let $G=\alpha X^e-\beta X^f\in I$ with $|e|,|f|\leq2$ and
$\alpha,\beta\in\mathbb Q$. In the target polynomial ring,
$\alpha t^{(Ae)_0}x^{(Ae)_1}y^{(Ae)_2}
=\beta t^{(Af)_0}x^{(Af)_1}y^{(Af)_2}$.
Equality of monomials implies either both coefficients vanish, or
$Ae=Af$ and $\alpha=\beta$. In either case $G=0$, by the injectivity
just proved.

If a set $S$ of such binomials generated $I$, each member of $S$
would lie in $I$ and hence be zero. Its generated ideal would then be
zero, contradicting $F\in I\setminus\{0\}$. Quadratic circuit relations
are a particular class of quadratic binomials in $I$, so they cannot
generate $I$ either.
\end{proof}

\section{Scope and formal verification}
The source~\cite{Source} asserts quadratic circuit generation in both
languages. Our configuration includes \emph{all} lattice points of an
explicit lattice polytope. The failure of this generation claim disproves
the conjecture; no conclusion about its Markov-basis clause is needed.

In \texttt{lean/Main.lean}, \texttt{Source} and \texttt{Target} are actual
Mathlib multivariate polynomial rings over $\mathbb Q$.
\texttt{toricMap} is an algebra homomorphism defined by polynomial
evaluation, and \texttt{toricIdeal} is its actual \texttt{RingHom.ker}.
The formal proof includes:
\begin{itemize}
\item \texttt{configuration\_is\_exact\_triangle} proves that the columns
come from exactly all the integer points in the triangle, using real
nonnegative barycentric coordinates.
\item \texttt{toricMap\_monomial} proves substitution for arbitrary
exponents and coefficients. The coordinate sum is proved equal to the
ordinary exponent sum.
\item The cubic is proved equal to $X_1X_2X_3-X_0^3$, nonzero, and in
the kernel. Every kernel binomial whose two terms have degree at most
two is proved zero.
\item \texttt{no\_quadratic\_binomial\_generating\_set} quantifies over
every set of such polynomials and rules out equality of its
\texttt{Ideal.span} with the kernel. No classification of arbitrary
quadratic polynomials or Markov bases is claimed.
\end{itemize}
The pinned project uses Lean 4.19.0 and Mathlib. Run in \texttt{lean/}:
\begin{verbatim}
lake build
lake env lean -DwarningAsError=true Main.lean
\end{verbatim}
The supplementary \texttt{verify.py} checks the lattice points, all
fifteen monomials of degree at most two, and the cubic with exact
arithmetic. Lean proves the universal ideal statement. Build logs, axiom
audits, and hashes are in \texttt{verification/}; the exact original
source bytes are in \texttt{SOURCE.md}.

\small
\begin{thebibliography}{9}
\bibitem{Source} The Last Math Competition, conjecture 00000003315,
\href{https://github.com/The-Last-Math-Competition/The-Last-Math-Competition/blob/main/conjectures/00000003315.md}{repository statement}.
\end{thebibliography}
\end{document}
